\clearpage
\section{Experiments and Results}

The Results chapter presents a deliberately structured progression from raw mechanical measurements to inferential analysis. Each plot represents a deliberate analytical choice: the figures are ordered to first establish the baseline dataset and spatial context, then characterise the physical dynamics of impact (deceleration, vibration, trajectory deviation), and finally assess whether the onboard IMU signatures alone can recover the impact angle without external MoCap ground truth. Comparative metrics between the Fixed Cage and Rotating Cage configurations are presented side-by-side throughout, enabling direct visual and statistical comparison of the bearing mechanism's effect.

To evaluate the operational and protective capabilities of the drone cage and to assess what indication of impact angle might be recoverable from onboard sensors, physical collision tests were conducted. This section details the experimental setup, the software pipeline used to capture the data, and the direct mechanical results, comparing the performance of the fixed cage against the freely rotating cage.

\subsection{Motion Capture Setup}

External state estimation is provided by a \textbf{6-camera OptiTrack Prime~22} system streaming at 120\,Hz via the Motive software platform (Mudfish network middleware). The Prime~22 offers sub-millimetre positional accuracy at short range, sufficient for the $\sim$5\,m$\times$5\,m capture volume used in these experiments.

The \texttt{motion\_capture\_tracking} ROS~2 node receives the multicast stream from Motive, and the \texttt{mocap\_px4\_bridge} node transforms the poses between coordinate frames before publishing them to \texttt{/fmu/in/vehicle\_visual\_odometry}.

\paragraph{Coordinate Frame Transformations}

Coordinate frame conventions differ between MoCap and PX4. ENU is a ground-fixed frame where the $X$~axis points East, $Y$~points North, and $Z$~points Up; the vehicle body frame has $X$~toward the front, $Z$~up, and $Y$~toward the left. PX4's internal state estimator operates in the \textbf{NED (North-East-Down)} frame, where $X$~points North, $Y$~points East, and $Z$~points Down; the corresponding vehicle body frame (\textbf{FRD}, Forward-Right-Down) has $X$~toward the front, $Z$~down, and $Y$~accordingly. \autoref{fig:coordinate_frames} shows both conventions side by side.\footnote{Software documentation available at: \url{https://docs.px4.io/main/en/tutorials/motion-capture}}

\begin{figure}[htbp]
    \centering
    \includegraphics[width=0.85\textwidth]{Figures/coordinate_frames_NEDvsENU.png}
    \caption{NED (left) and ENU (right) coordinate frame conventions. Adapted from PX4 documentation~\footnote{Software documentation available at: \url{https://docs.px4.io/main/en/tutorials/motion-capture}}.}
    \label{fig:coordinate_frames}
\end{figure}

The \texttt{mocap\_px4\_bridge} performs this transform in two steps: first, a static rotation maps ENU to NED ($X_{\text{NED}} = Y_{\text{ENU}}$, $Y_{\text{NED}} = X_{\text{ENU}}$, $Z_{\text{NED}} = -Z_{\text{ENU}}$); second, the NED pose is rotated into the FRD body frame using the vehicle's current attitude quaternion. The resulting \texttt{VehicleVisualOdometry} message is published at 120\,Hz to \texttt{/fmu/in/vehicle\_visual\_odometry}.

The EKF2 estimator fuses this external vision with the onboard 250\,Hz IMU data. During steady, unoccluded flight, this fusion produces velocity estimates that align closely with the MoCap-derived ground truth, as shown in Figure~\ref{fig:ekf_validation} for representative flights from both cage configurations.

\paragraph{EKF Validation Under MoCap Occlusion}

During the exact millisecond of physical impact, the protective cage and column structure physically block the infrared camera's line-of-sight to the drone's reflective markers. This transient marker occlusion causes the MoCap tracking data to glitch. When velocity and acceleration are derived from position by numerical differentiation, even a single-frame position glitch produces a massive, non-physical spike in the resulting kinematic profile. The Savitzky-Golay differentiation filter cannot remove this artefact---it amplifies any discontinuity in the underlying data. Figure~\ref{fig:mocap_dropout} illustrates this failure: the MoCap-derived velocity and acceleration traces exhibit violent, unphysical excursions at the impact instant ($t = 0$), while the EKF2 odometry remains smooth and continuous throughout the event.

The EKF2 does not rely solely on the external vision stream. It continuously fuses the onboard 250\,Hz IMU accelerometer and gyroscope data into its state estimate. When the MoCap feed momentarily degrades, the inertial propagation bridges the gap without introducing noticeable drift. The smooth, physically plausible EKF velocity profile through the collision window---visible in Figure~\ref{fig:mocap_dropout}---confirms that the filter successfully compensates for the optical dropout.

\begin{figure}[htbp]
    \centering
    \includegraphics[width=0.95\textwidth]{plots/Fixed Cage Kinetic Profile (Raw).png}
    \caption{Figure A: MoCap dropout during impact. The Savitzky-Golay filtered MoCap velocity (blue, top panel) and acceleration (green, bottom panel) exhibit severe non-physical spikes at $t = 0$, caused by marker occlusion at the instant of column contact. The EKF2 odometry (red) remains smooth and continuous, bridging the dropout via inertial propagation.}
    \label{fig:mocap_dropout}
\end{figure}

\begin{figure}[htbp]
    \centering
    \includegraphics[width=0.95\textwidth]{plots/EKF Dual Comparison Fixed vs Rotating.png}
    \caption{Figure B: EKF vs.\ MoCap velocity validation during clean, unoccluded flight for representative Fixed Cage (left) and Rotating Cage (right) flights. The close alignment between the two traces confirms that the EKF2 accurately mirrors the MoCap ground truth under normal tracking conditions.}
    \label{fig:ekf_validation}
\end{figure}

Based on this validated alignment, EKF odometry is used as the primary source of velocity and acceleration data for all kinematic calculations through the collision window. The raw MoCap position data is retained solely for geometric reference (impact angle computation, clearance measurement, and trajectory visualization)---calculations that do not require numerical differentiation and are therefore insensitive to transient position glitches.

\begin{figure}[htbp]
    \centering
    \includegraphics[width=0.85\textwidth]{plots/Impact Angle Distribution by Cage Configuration.png}
    \caption{Distribution of measured impact angles for both cage configurations. The rotating cage sample is shifted toward the lower end of the $45^\circ$ approach condition (mean~$\pm$~SD: $40.6 \pm 11.5^\circ$) compared to the fixed cage ($46.7 \pm 13.0^\circ$), reflecting variance in collision geometry across passes.}
    \label{fig:impact_angle_distribution}
\end{figure}

\subsection{Experimental Setup}
\subsubsection{Experiment Environment and Collision Target}
The physical experiments were conducted at the REALLAB at the IT University of Copenhagen (ITU) within an OptiTrack Motion Capture (MoCap) environment. A vertically oriented cylindrical cardboard column, internally reinforced with an aluminum profile, was anchored to the floor and ceiling to serve as the static collision target. The column oscillates slightly when subjected to manual force, but no additional structural supports were added in order to prevent any occlusion of the MoCap cameras, preserving their function as the primary reference point for the drone's spatial positioning. Analysis of the MoCap data confirmed that obstacle displacement during impacts was below the system's millimetre tracking resolution. Therefore, energy absorption by the environment is considered negligible.

\subsubsection{Experiment Control and Data Collection}
The experimental platform relies on an integrated software architecture designed to bridge abstracted trajectory generation with real-time hardware constraints. The low-level flight stack utilizes PX4 Autopilot firmware on a NuttX real-time operating system (RTOS), with internal state estimation handled strictly by an Extended Kalman Filter (EKF2).\footnote{Software documentation available at: \url{https://docs.px4.io/main/en/flight_controller/pixhawk6c}} \textbf{L1 Node} --- the EKF2 is the primary source of velocity state data for the entire analysis pipeline, fused from the onboard IMU data logged at 250\,Hz (the rate at which the EKF2 pipeline ingests the raw accelerometer and gyroscope data) and the 120\,Hz external MoCap vision odometry stream.

High-level position commanding is managed within a ROS~2 abstraction layer running on an external offboard computing node. Real-time topic translation between the ROS~2 node graph and the internal PX4 uORB messaging system is handled via a micro-XRCE-DDS agent.

To orchestrate the experimental runs, a custom software component named \texttt{flight\_director.py} was developed. This script functions as an automated mission execution environment, systematically enforcing geofence safety parameters. The flight director automates the data collection pipeline by initializing synchronized telemetry recordings, saving each distinct flight loop into an independent \texttt{.mcap} file. These files were subsequently parsed using a custom Python database pipeline to calculate the physical variables presented in the following results.

The complete experimental pipeline, from hardware components through the ROS~2 software stack and MoCap integration to telemetry storage, is illustrated in \autoref{fig:experiments_pipeline}.

\subsubsection{Data Processing for Collision Analysis}
\label{sec:mocap_ekf2}

During the exact millisecond of physical impact, the protective cage structures transiently occlude the drone's reflective markers. This occlusion frequently caused the MoCap tracking publish rate to drop below the critical 30\,Hz failsafe threshold, occasionally plummeting to $\sim$10\,Hz. When relying purely on raw spatial differentiation, these temporal gaps produced massive, artificial high-frequency spikes in the calculated kinematic profiles, even when applying Savitzky-Golay filtering.

To resolve this without sacrificing data fidelity, the analysis pipeline bypasses raw spatial differentiation for velocity and acceleration entirely. Instead, velocity states are extracted directly from the internal PX4 Extended Kalman Filter (\texttt{vehicle\_odometry.velocity}). Because the EKF2 continuously fuses the onboard 250\,Hz IMU data, it bridges external vision dropouts using inertial propagation, providing a smooth and physically truthful velocity profile throughout the collision event.

\subsubsection{Pre-Computed Path and Waypoint Configuration}

To reiterate, this physical collision framework represents the actual experimental flight tests conducted, from which all the data in this thesis stems. The system uses a pre-computed ROS~2 waypoint trajectory to ensure every pass is mechanically identical. The exact spatial path, including the approach, impact, and recovery segments, is detailed in Figure~\ref{fig:mission_geometry}.

\begin{figure}[h!]
    \centering
    \includegraphics[width=\textwidth]{plots/45° Column Collision Loop Full Mission Geometry.png}
    \caption{Full mission geometry for the $45^\circ$ collision experiment. The dotted red line traces the closed-loop waypoint path. Waypoints correspond to the state machine: \texttt{WP\_stage} initiates the forward approach; Exp.\ Start-point (\texttt{WP1}) marks the start of the high-momentum sweep leg; Exp.\ End-point (\texttt{WP2}) marks the exit; and Recovery (\texttt{WP3}) is the hold position. \textit{An additional \texttt{WP0} (Takeoff) exists off-diagram at the drone's physical starting location on the floor, utilized only once at the start of the flight. The column obstacle is shown at the centre. Note that the `ideal' line between Exp.\ Start-point (\texttt{WP1}) and Exp.\ End-point (\texttt{WP2}) dictates the trajectory, and Exp.\ End-point (\texttt{WP2}) possesses a 15~cm diameter acceptance leeway to register the waypoint as reached.}}
    \label{fig:mission_geometry}
\end{figure}

% ======================================================================
\subsection{Experimental Kinematics and Spatial Recovery}
\label{sec:deviation_from_norm}
% ======================================================================

This section reports the macro-spatial kinematics of the collisions. As detailed in Table~\ref{tab:dataset_baseline}, the analysis is based on 179 executed passes, yielding 168 recorded passes and 127 verified impacts across both configurations.

\begin{table}[h!]
\centering
\caption{Dataset baseline: total flight passes, verified structural impacts, and the isolation criteria used to distinguish genuine collisions from near-miss transits.}
\label{tab:dataset_baseline}
\resizebox{\textwidth}{!}{%
\begin{tabular}{lcc}
\toprule
\textbf{} & \textbf{Count} & \textbf{Notes} \\
\midrule
Total flight passes executed & 179 & Across all experimental sessions \\
\addlinespace
Verified structural impacts & 127 & Confirmed via telemetry inspection + cage-radius verification \\
\addlinespace
\hspace{0.5em}-- Fixed Cage & 67 & \\
\hspace{0.5em}-- Rotating Cage & 60 & \\
\addlinespace
Isolation criterion & -- & Physical contact confirmed when $d_{\text{centre}} < r_{\text{column}} + r_{\text{cage}} = 0.224$~m \\
\addlinespace
Battery drain rate (\%\,min$^{-1}$) & 21.3 & Both configurations; full unsliced flight sessions \\
\bottomrule
\end{tabular}}%
\end{table}

\subsubsection{Trajectory Visualisation and Mission Outcomes}
\label{sec:horizontal_kinematics}

To evaluate the drone's horizontal flight dynamics, this analysis focuses on the spatial divergence between the pre-computed commanded path and the actual trajectory executed. Visualizing this path deviation in Figure~\ref{fig:horizontal_kinematics} displays the physical trajectories executed during the impacts. A breakdown of the specific mission outcomes (impacts vs.\ no impacts across angles) is provided in Table~\ref{tab:mission_outcomes}.

\begin{figure}[h!]
    \centering
    \includegraphics[width=\textwidth]{plots/Experiment 2D horizontal visualization.png}
    \caption{2D horizontal visualization of the experiment flight paths. The plot compares commanded versus executed trajectories for both fixed and rotating cage configurations across the collision sweep leg, showing spatial deviation, impact coordinates, and recovery overshoot required to stabilize lateral drift.}
    \label{fig:horizontal_kinematics}
\end{figure}

\begin{table}[h!]
    \centering
    % Auto-generated by plot_mission_outcome_table()
% Requires: \usepackage{booktabs}
\begin{tabular}{lrrrrrrr}
\toprule
 & \rotatebox{90}{\textsf{N}} & \rotatebox{90}{\textsf{No Impact}} & \rotatebox{90}{\textsf{{<\,}30°}} & \rotatebox{90}{\textsf{30–40°}} & \rotatebox{90}{\textsf{40–50°}} & \rotatebox{90}{\textsf{50–60°}} & \rotatebox{90}{\textsf{60–90°}} \\
\midrule
\multicolumn{8}{l}{\textbf{Rotating Cage}} \\
\midrule
\quad 45° & 50 & 2 & 9 & 21 & 13 & 5 & 0 \\
\quad 75° & 31 & 19 & 0 & 2 & 2 & 3 & 5 \\
\quad \textbf{All} & 81 & 21 & 9 & 23 & 15 & 8 & 5 \\
\addlinespace
\multicolumn{8}{l}{\textbf{Fixed Cage}} \\
\midrule
\quad 45° & 56 & 5 & 3 & 18 & 13 & 10 & 7 \\
\quad 75° & 31 & 15 & 2 & 0 & 4 & 5 & 5 \\
\quad \textbf{All} & 87 & 20 & 5 & 18 & 17 & 15 & 12 \\
\bottomrule
\end{tabular}
%% Data: experiments_summary.db → flights_summary (81× Rotating, 87× Fixed flights)

    \caption{Mission outcome summary across all 168 recorded passes. Flights are categorised by cage configuration, target approach angle ($45^\circ$ and $75^\circ$), and the measured impact angle at the collision instant. The ``No Impact'' column counts passes where the drone cleared the column without physical contact.}
    \label{tab:mission_outcomes}
\end{table}

\subsubsection{Spatial Trajectory and Recovery Metrics}
\label{sec:clearance_sdld}

The spatial consequences of the physical impacts are quantified by assessing the clearance, the worst-case deviation, and the total recovery area. These values are summarized in Table~\ref{tab:spatial_metrics}.

\textbf{Obstacle clearance} represents the closest physical approach. The rotating cage exhibits a larger mean negative clearance, meaning it incurs deeper into the column boundary (Rotating: $\mu = -0.637 \pm 0.172$~m; Fixed: $\mu = -0.470 \pm 0.165$~m).

\textbf{Maximum Trajectory Deviation} ($d_{\text{max}}$) is the peak perpendicular distance between the executed MoCap track and the commanded `ideal' reference line. This ideal line is drawn exactly from the coordinate where the drone was commanded forward (\texttt{WP1}) to the exact target coordinate it must reach (\texttt{WP2}, which includes a 15~cm precision leeway). This metric captures the single worst-case excursion during the obstacle transit. According to Table~\ref{tab:spatial_metrics}, the rotating cage recorded a maximum trajectory deviation of $131.0 \pm 32.1$~mm, while the fixed cage recorded $118.4 \pm 30.4$~mm.

\textbf{Spatial Integral of Absolute Error} (SIAE) measures the total spatial area footprint of the recovery phase. Table~\ref{tab:spatial_metrics} shows the spatial area of the recovery phase was $120.4 \pm 32.2 \times 10^{-3}~\text{m}^2$ for the rotating cage and $117.4 \pm 33.2 \times 10^{-3}~\text{m}^2$ for the fixed cage.

Finally, the \textbf{Path Spread (RMSLD)} quantifies how tightly the drone tracks the ideal line over the entirety of the obstacle passage. To visualize this overall consistency across all flight tests, the overlaid paths of all flights are shown in Figure~\ref{fig:path_heatmap}. As extracted from the dataset, the path spread is recorded at $40.3 \pm 9.8$~mm for the rotating cage and $37.4 \pm 11.0$~mm for the fixed cage.

\begin{figure}[htbp]
    \centering
    \includegraphics[width=\textwidth]{plot_16_path_heatmap.png}
    \caption{Path overlay of all impact flights for both cage configurations. Each flight trajectory is plotted at low opacity, showing the aggregate spatial footprint of the collision sweep for Rotating Cage (left, blue) and Fixed Cage (right, red). The dashed line marks the nominal commanded path.}
    \label{fig:path_heatmap}
\end{figure}

Based on the data reported above, the spatial recovery metrics suggest that the rotating cage's bearing mechanism does not expand the final recovery footprint required by the standard flight controller settings. This pattern is further evaluated by comparing the worst-case excursions and recovery consistency across configurations.

The two complementary metrics---maximum deviation and SIAE---confirm that a brief large excursion and a prolonged small offset can produce similar SIAE values. The rotating cage exhibited a slightly wider maximum trajectory deviation than the fixed cage. A Mann-Whitney U test confirms this difference is statistically significant ($U = 2484$, $p = 0.022$), though the effect size is small (Cohen's $d = 0.41$). However, the total spatial area required by the PID controller to arrest the drift and return to the setpoint (SIAE) remains nearly identical between both configurations ($U = 2009$, $p = 0.998$, not significant).

\begin{table}[htbp]
\centering
\caption{Spatial and Clearance Metrics. Summary of key comparative spatial metrics between cage configurations. All values are mean~$\pm$~standard deviation.}
\label{tab:spatial_metrics}
\resizebox{\textwidth}{!}{%
\begin{tabular}{lccc}
\toprule
\textbf{Metric} & \textbf{Rotating Cage} & \textbf{Fixed Cage} & \textbf{Statistical Test} \\
\midrule
Obstacle Clearance (m) & $-0.637 \pm 0.172$ & $-0.470 \pm 0.165$ & --- \\
\addlinespace
Max. Trajectory Deviation (mm) & $131.0 \pm 32.1$ & $118.4 \pm 30.4$ & $U=2484$, $p=0.022$, $d=0.41$ \\
\addlinespace
SIAE Recovery Area (\texttimes 10$^{-3}$\,m$^{2}$) & $120.4 \pm 32.2$ & $117.4 \pm 33.2$ & $U=2009$, $p=0.998$\,(n.s.) \\
\addlinespace
Path Spread RMSLD (mm) & $40.3 \pm 9.8$ & $37.4 \pm 11.0$ & --- \\
\bottomrule
\end{tabular}}%
\end{table}

% ======================================================================
\subsection{Impact Acceleration and Vibrations}
\label{sec:impact_acceleration_vibrations}
% ======================================================================

To understand the physical forces at contact, we evaluate the IMU acceleration and gyroscopic data. During these extreme stress tests---consisting of repeated, high-throttle crash recoveries---the platform averaged a heavy continuous battery drain of 21.3\% per minute. Because such voltage drops can affect motor authority, it was necessary to confirm that battery depletion did not secretly skew the deceleration data. The covariance with the starting battery percentage is summarised in Table~\ref{tab:deceleration_battery}.\label{sec:efficiency_baseline}

\begin{table}[h!]
\centering
\caption{Impact Deceleration and Battery Covariance. Summary of impact deceleration and its covariance with starting battery percentage. The near-zero Huber slopes confirm that battery state-of-charge did not confound the deceleration measurements. Note: The relative difference can be expressed either as a 37.2\% reduction from the fixed baseline, or equivalently, the fixed cage experiencing a 59\% higher deceleration than the rotating cage.}
\label{tab:deceleration_battery}
\begin{tabular}{lcc}
\toprule
\textbf{Metric} & \textbf{Fixed Cage} & \textbf{Rotating Cage} \\
\midrule
Mean Max.\ Deceleration ($\text{m/s}^2$) & $-1.64 \pm 0.64$ & $-1.03 \pm 0.50$ \\
Relative Difference & Baseline & $-37.2\%$ reduction \\
Battery Covariance (Huber Slope) & $-0.005$ & $0.000$ \\
\bottomrule
\end{tabular}
\end{table}

As detailed in Table~\ref{tab:deceleration_battery}, the rotating cage recorded a mean maximum deceleration of $-1.03 \pm 0.50$~m/s$^2$, while the fixed cage recorded $-1.64 \pm 0.64$~m/s$^2$. To clarify the relative performance differences, these values are expressed as absolute magnitudes (representing the raw physical force transferred to the airframe).

When evaluating how much the fixed cage amplifies the impact forces relative to the rotating cage, the rotating cage serves as the baseline ($D_{\text{rot}} = 1.03$):
\[ \text{Relative Increase} = \frac{|D_{\text{fix}}| - |D_{\text{rot}}|}{|D_{\text{rot}}|} = \frac{1.64 - 1.03}{1.03} \approx 0.592 \]
This demonstrates a 59.2\% higher maximum deceleration for the fixed cage.

Conversely, when evaluating the mitigating effect of the bearing mechanism, the standard fixed cage serves as the baseline ($D_{\text{fix}} = 1.64$):
\[ \text{Relative Reduction} = \frac{|D_{\text{fix}}| - |D_{\text{rot}}|}{|D_{\text{fix}}|} = \frac{1.64 - 1.03}{1.64} \approx 0.372 \]
This demonstrates a 37.2\% reduction in impact deceleration achieved by the rotating cage.

An independent samples t-test (Welch) confirms this difference is statistically significant ($t(124.7) = 6.00$, $p < 0.001$, Cohen's $d = 1.05$). This strong statistical significance points to a regular, normalized data distribution, confirming the high reliability and relevancy of the obtained deceleration values.

\paragraph{IMU Time-Series: Individual and Aggregated Signatures}

Figure~\ref{fig:imu_dynamics_pair} presents the acceleration and gyroscopic traces. The upper panel shows representative single-flight traces. The lower panel aggregates all 127 flights. In the aggregated data, the overlaid individual traces reveal a visually narrower spread of linear acceleration peaks for the rotating cage, alongside a shorter median gyroscopic peak duration compared to the fixed cage.

\begin{figure}[htbp]
    \centering
    \begin{subfigure}{\textwidth}
        \centering
        \includegraphics[width=\textwidth]{plots/Combined IMU Collision Dynamics.png}
        \caption{Representative single-flight traces for Rotating Cage (left) and Fixed Cage (right).}
        \label{fig:combined_imu_dynamics}
    \end{subfigure}
    \vspace{1em}
    \begin{subfigure}{\textwidth}
        \centering
        \includegraphics[width=\textwidth]{plots/Aggregated IMU Collision Dynamics Across All Flights.png}
        \caption{Aggregate overlay of all 127 flights with median profiles.}
        \label{fig:aggregated_imu_dynamics}
    \end{subfigure}
    \caption{IMU collision dynamics at two granularities: (a) representative individual flights reveal the per-event shock and decay signature; (b) aggregation across all 127 flights confirms the pattern is systematic. Together they establish that the rotating cage consistently attenuates both the magnitude and the variability of impact-induced vibration.}
    \label{fig:imu_dynamics_pair}
\end{figure}

\paragraph{Axis-Resolved IMU Component Analysis}

Disaggregating by axis, Figure~\ref{fig:raw_imu_xyz} shows the X, Y, Z collision window. The Y axis (lateral) records peaks below $\pm 4~\text{m/s}^2$ for the rotating cage, compared to peaks exceeding $\pm 6~\text{m/s}^2$ for the fixed cage.

\begin{figure}[h!]
    \centering
    \includegraphics[width=\textwidth]{plots/Combined IMU XYZ Components.png}
    \caption{Raw acceleration traces for each axis (X, Y, Z) during the collision window. The $Y$~axis (lateral) carries the dominant collision signature in both configurations; the rotating cage reduces peak amplitude by roughly one-third on this channel.}
    \label{fig:raw_imu_xyz}
\end{figure}

The box plots in Figure~\ref{fig:xyz_oscillation_variance} summarise the post-impact stabilisation tail ($t_{\text{impact}} + 0.2\,\text{s}$ to $+3.0\,\text{s}$) across all flights. During steady flight, the rotating cage exhibits a lower variance spread in both the X and Y axes. During impact, the rotating cage records a lower spread in the Y axis, but registers a larger interquartile spread in the Z axis.

\begin{figure}[htbp]
    \centering
    \includegraphics[width=\textwidth]{plots/IMU Vibration Spread by Cage Configuration.png}
    \caption{Box plots of per-axis variance over the post-impact stabilisation tail ($t_{\text{impact}} + 0.2\,\text{s}$ to $+3.0\,\text{s}$). The rotating cage's narrower interquartile ranges confirm that the bearing reduces both peak shock and the duration of residual vibration.}
    \label{fig:xyz_oscillation_variance}
\end{figure}

To provide a holistic view of the drone's subsystem performance, Table~\ref{tab:master_comparison} aggregates the key metrics across kinematics, motor output, attitude control, and vibration. This table highlights several noteworthy differences, specifically that the rotating cage records substantially lower yaw-rate tracking error ($-85\%$) and reduced allocator saturation ($-96\%$) compared to the fixed configuration.

\paragraph{EKF State Estimation Fidelity}

Because the MoCap markers can become occluded during physical contact with the column, it was necessary to verify if the flight controller's internal state estimation could reliably bridge these visual dropouts. Figure~\ref{fig:ekf_validation} plots the EKF-estimated velocity alongside the MoCap-derived velocity over the collision window for representative flights. The two traces track each other closely through the entire collision window. This agreement confirms that the EKF2's inertial propagation reliably bridges the transient MoCap dropout without introducing noticeable drift, justifying the use of EKF telemetry for the collision analysis.

\medskip
The preceding kinematic and vibrational analyses establish that the rotating cage fundamentally alters the physical distribution of collision energy. However, for a drone operating without external motion-tracking, this mechanical advantage is only practically useful if the onboard sensors retain enough information to characterise the impact. As indicated by the IMU feature correlation analysis (see \autoref{fig:imu_correlation}), the relationship between impact angle and inertial response is highly complex and dispersed across multiple axes. This dispersion suggests that simple linear thresholds are insufficient, necessitating a multivariate, non-linear approach to successfully decode the collision geometry from the raw telemetry.

\begin{table}[h!]
\centering
\caption{Comprehensive Performance Summary. Comprehensive comparison of all key performance metrics between cage configurations. Values are mean~$\pm$~standard deviation across all 127 verified impact flights (60 Rotating Cage, 67 Fixed Cage). Arrow direction ($\leftarrow$ or $\rightarrow$) indicates which configuration achieves the favourable outcome for each metric; the $\Delta$ column shows the percentage difference relative to the Fixed Cage baseline.}
\label{tab:master_comparison}
\footnotesize
\setlength{\tabcolsep}{3pt}
\begin{tabular}{lcccr}
\toprule
\textbf{Category / Metric} & \textbf{Rotating} & \textbf{Fixed} & \textbf{$\Delta$ (\%)} & \textbf{Better} \\
\midrule
\multicolumn{5}{l}{\textbf{Kinematics}} \\
~~Peak Decel. Z (m/s$^{2}$) & $10.0\pm4.7$ & $6.57\pm4.00$ & $+52.5$ & $\rightarrow$ Fixed \\
~~Max Traj. Deviation (mm) & $131.0\pm32.1$ & $118.4\pm30.4$ & $+10.7$ & $\rightarrow$ Fixed \\
~~Integ. Rot. Energy (rad) & $0.16\pm0.05$ & $0.24\pm0.07$ & $-33.7$ & $\leftarrow$ Rotating \\
~~Recovery Area (m$^{2}$) & $0.12\pm0.03$ & $0.12\pm0.03$ & $+2.6$ & $\rightarrow$ Fixed \\
~~Path Spread RMSLD (mm) & $40.3\pm9.8$ & $37.4\pm11.0$ & $+7.8$ & $\rightarrow$ Fixed \\
~~Experiment Duration (s) & $7.94\pm1.26$ & $7.88\pm1.76$ & $+0.7$ & $\rightarrow$ Fixed \\
\midrule
\multicolumn{5}{l}{\textbf{Motors \& Energy}} \\
~~Max Actuator Output (\%) & $99.3\pm1.5$ & $99.9\pm0.0$ & $-0.6$ & $\leftarrow$ Rotating \\
\midrule
\multicolumn{5}{l}{\textbf{Attitude Control}} \\
~~Roll Rate RMS (rad/s) & $0.22\pm0.05$ & $0.16\pm0.05$ & $+36.2$ & $\rightarrow$ Fixed \\
~~Pitch Rate RMS (rad/s) & $0.29\pm0.12$ & $0.25\pm0.11$ & $+13.4$ & $\rightarrow$ Fixed \\
~~Yaw Rate RMS (rad/s) & $0.09\pm0.03$ & $0.57\pm0.15$ & $-85.0$ & $\leftarrow$ Rotating \\
\midrule
\multicolumn{5}{l}{\textbf{Control Allocation}} \\
~~Alloc. Sat. Duration (s) & $0.01\pm0.02$ & $0.13\pm0.10$ & $-96.2$ & $\leftarrow$ Rotating \\
~~Max Unalloc. Torque (N\,m) & $0.01\pm0.04$ & $0.01\pm0.02$ & $-2.3$ & $\leftarrow$ Rotating \\
~~Thrust Setpt. Achieved (\%) & $97.3\pm6.9$ & $85.1\pm13.2$ & $+14.4$ & $\leftarrow$ Rotating \\
\midrule
\multicolumn{5}{l}{\textbf{IMU Vibration}} \\
~~Accel X Impact (g) & $0.28\pm0.11$ & $0.29\pm0.13$ & $-3.6$ & $\leftarrow$ Rotating \\
~~Accel Y Impact (g) & $0.26\pm0.08$ & $0.37\pm0.16$ & $-28.6$ & $\leftarrow$ Rotating \\
~~Accel Z Impact (g) & $0.24\pm0.09$ & $0.18\pm0.09$ & $+32.3$ & $\rightarrow$ Fixed \\
~~Accel X Regular (g) & $0.01\pm0.01$ & $0.05\pm0.07$ & $-69.5$ & $\leftarrow$ Rotating \\
~~Accel Y Regular (g) & $0.02\pm0.01$ & $0.10\pm0.11$ & $-77.1$ & $\leftarrow$ Rotating \\
~~Accel Z Regular (g) & $0.05\pm0.01$ & $0.04\pm0.03$ & $+10.8$ & $\rightarrow$ Fixed \\
\midrule
\bottomrule
\end{tabular}
\end{table}

% ======================================================================
\subsection{Intermediate Mechanical Discussion}
\label{sec:intermediate_discussion}
% ======================================================================

Based on the data reported in Sections~\ref{sec:deviation_from_norm} and~\ref{sec:impact_acceleration_vibrations}, several mechanical hypotheses can be formulated regarding the bearing mechanism's effect on collision energy.

First, the trajectory recovery area (SIAE) remained nearly identical between configurations despite a slightly wider initial deviation for the rotating cage. This indicates that the spatial cost of the rotating mechanism likely does not expand the final recovery footprint required by the standard flight controller settings.

Second, the gyroscopic and acceleration data suggest that the rotating cage alters the physical distribution of collision energy. The bearing mechanism seemingly does not uniformly dampen all vibrations. While it mitigates lateral (Y-axis) shock, the increased variance in the Z-axis (Figure~\ref{fig:xyz_oscillation_variance}) indicates the cage likely redirects a portion of the lateral collision momentum into vertical airframe excitation.

Because the rotating cage appears to scatter the impact energy across different axes rather than absorbing it linearly, it is hypothesized that simple threshold models will be mathematically insufficient to determine the collision geometry. This dispersion effect necessitates the multivariate, non-linear machine learning approach applied in the following section.

\clearpage

% ======================================================================
\subsection{Sensor-Based Impact Angle Inference}
\label{sec:result_analysis}
% ======================================================================

We have established that the crash generates complex, multi-axis forces and that the rotating cage fundamentally changes their distribution. The practical question now becomes: can a computer use those onboard IMU forces alone --- without any external MoCap ground truth --- to determine the angle at which the drone struck the column?

\begin{figure}[htbp]
    \centering
    \includegraphics[width=\textwidth]{plots/consolidated_feature_correlation.png}
    \caption{Consolidated feature correlation slopegraph: Pearson~$r$ between each of the 26 IMU aggregate metrics and \texttt{impact\_angle}, for Rotating Cage (left heatmap) and Fixed Cage (right heatmap). Splines connecting the two columns trace how each feature's rank shifts between cage types --- red splines indicate large rank shifts ($\geq$4 positions); grey splines indicate stable ranks. Significance stars: *** $p<0.001$, ** $p<0.01$, * $p<0.05$.}
    \label{fig:imu_correlation}
\end{figure}

This dispersion is further visualised through parallel coordinates (\autoref{fig:consolidated_parallel_coordinates}), which display the overlapping, multivariate trajectories of the top 10 IMU features ordered by descending $|\text{Pearson}~r|$ (computed on Rotating Cage data). Each polyline represents one flight, coloured by impact-angle bin. The dense entanglement of these polylines across angle bins visually demonstrates why simple linear thresholds fail --- no single axis cleanly separates the angle groups. When individual flights are traced across the top features, the polyline bundles for different angle bins intertwine heavily with no clean separation surface. This chaotic, overlapping structure is the direct visual evidence that the angle signal is distributed across multiple channels and cannot be recovered by a single linear threshold, mathematically justifying the transition to the Random Forest architecture that follows.

\begin{figure}[htbp]
    \centering
    \includegraphics[width=\textwidth]{plots/consolidated_parallel_coordinates.png}
    \caption{Parallel coordinates of the top IMU features against impact angle. Rotating Cage (top) and Fixed Cage (bottom). Polylines are coloured by angle bin.}
    \label{fig:consolidated_parallel_coordinates}
\end{figure}

The correlation slopegraph (\autoref{fig:imu_correlation}) reveals that no individual IMU feature achieves a Pearson correlation above $|r| \approx 0.6$ with impact angle in either cage configuration, and that feature ranks shift substantially between the two mechanical regimes — confirming that the angle signal is distributed across multiple channels and cannot be recovered by a single linear threshold.

\subsubsection{Angle Prediction via Random Forest}
\label{sec:angle_prediction}

To map these complex, dispersed vibrational signatures to a specific impact angle, a Random Forest regressor was employed. Random Forest is an ensemble machine learning algorithm that constructs multiple decision trees during training and outputs the mean prediction of those individual trees. It is particularly well-suited for this application because it natively models non-linear relationships and complex interactions between multiple variables---such as the coupled gyroscopic and acceleration forces generated during a multi-axis collision. Furthermore, unlike strict linear models, it does not require the underlying physical data to follow a formal mathematical distribution, making it highly effective at parsing chaotic, high-frequency crash telemetry.

Model performance is evaluated using three standard regression metrics to provide a holistic view of predictive accuracy. \textbf{Mean Absolute Error (MAE)} represents the average absolute difference between the predicted and actual impact angles, providing a direct, easily interpretable measure of typical angular error in degrees. \textbf{Root Mean Square Error (RMSE)} similarly measures prediction error but mathematically penalizes larger deviations more heavily; a substantial gap between MAE and RMSE indicates the presence of severe individual mispredictions. Finally, the \textbf{Coefficient of Determination ($R^2$)} quantifies the proportion of variance in the true impact angle that the IMU features successfully explain, serving as a normalized indicator of the model's overall predictive power.

While the preceding sections characterise the mechanical consequences of collision, they rely on post-processed MoCap ground truth. A practical drone operating in a confined space does not have access to an external tracking system --- it must infer collision properties from onboard sensors alone. To evaluate whether impact angle can be estimated from the vibration and inertial signatures recorded by the onboard IMU, a Random Forest regressor was trained on the 14 IMU-derived features described in the methodology (\autoref{sec:mocap_ekf2}). The model was evaluated separately for each cage configuration using stratified 5-fold cross-validation, enabling direct comparison of how well the vibration-to-angle mapping survives the mechanical differences between fixed and rotating designs.

\subsubsection{Five-Fold Cross-Validation: Predicted vs.\ Actual Impact Angle}

The 5-fold cross-validation procedure splits the impact flights for each cage configuration into five subgroups, trains the model on four, and evaluates on the held-out fifth in rotation. This ensures every flight contributes to testing exactly once, providing an unbiased estimate of generalisation performance on unseen data.

\begin{figure}[htbp]
    \centering
    \includegraphics[width=\textwidth]{plots/consolidated_actual_vs_predicted.png}
    \caption{Consolidated five-fold cross-validation: predicted vs.\ actual impact angle for both cage configurations. Fixed Cage (left) and Rotating Cage (right) are shown side by side with identical axis limits, aspect ratio, and tick spacing for direct visual comparison. Each marker colour denotes one cross-validation fold; the black dashed line marks the $y=x$ identity; the grey band shows $\pm5^\circ$ tolerance.}
    \label{fig:rf_cv_consolidated}
\end{figure}

Both configurations produced predictive performance substantially above chance, though with distinct error profiles (Table~\ref{tab:rf_cv_metrics}, \autoref{fig:rf_cv_consolidated}). The Rotating Cage achieved lower absolute error (MAE = 4.53$^\circ$, RMSE = 6.06$^\circ$) than the Fixed Cage (MAE = 4.87$^\circ$, RMSE = 6.71$^\circ$), indicating its predictions were physically closer to the true impact angles. The Fixed Cage, however, returned a slightly higher $R^2$ (0.731 vs.~0.719), reflecting stronger overall correlation between its predicted and actual trends despite marginally larger absolute deviations. This apparent paradox --- higher $R^2$ but higher MAE/RMSE --- arises because $R^2$ measures the proportion of variance explained, not the absolute magnitude of error, and the Fixed Cage angle distribution spans a wider dynamic range.

\begin{table}[htbp]
\centering
\caption{Cross-validated Random Forest performance metrics for impact-angle prediction by cage configuration. All metrics aggregated across five folds; Huber baseline uses the single best-correlated IMU feature.}
\label{tab:rf_cv_metrics}
\resizebox{\textwidth}{!}{%
\begin{tabular}{lccc}
\toprule
\textbf{Metric} & \textbf{Fixed Cage} & \textbf{Rotating Cage} & \textbf{Direction} \\
\midrule
CV $R^2$ & 0.731 & 0.719 & Higher is better \\
\addlinespace
MAE ($^\circ$) & 4.87 & 4.53 & Lower is better \\
\addlinespace
RMSE ($^\circ$) & 6.71 & 6.06 & Lower is better \\
\addlinespace
OOB $R^2$ & 0.706 & 0.693 & Higher is better \\
\addlinespace
Huber baseline $R^2$ & 0.348 & 0.004 & Higher is better \\
\addlinespace
Features selected & 14 & 14 & --- \\
\bottomrule
\end{tabular}}%
\end{table}

The residual plots (\autoref{fig:rf_residuals_consolidated}) reveal a diagnostic pattern consistent across both cage types. In both configurations, but particularly visible in the Rotating Cage, the residuals exhibit a discernible curve running from negative to positive rather than a fully random scatter around zero. This structured pattern indicates that the underlying relationship between IMU signatures and impact angle possesses a non-linear component that the current Random Forest configuration does not completely capture --- the model leaves a systematic bias that a deeper or differently tuned architecture could potentially absorb.

\begin{figure}[htbp]
    \centering
    \includegraphics[width=\textwidth]{plots/consolidated_residuals.png}
    \caption{Consolidated residual plots for both cage configurations. Rotating Cage (left) and Fixed Cage (right). The black dashed horizontal line marks zero error; residuals are coloured by cross-validation fold.}
    \label{fig:rf_residuals_consolidated}
\end{figure}

\subsubsection{Learning Curves and Overfitting}

\begin{figure}[htbp]
    \centering
    \includegraphics[width=0.85\textwidth]{plots/consolidated_learning_curves.png}
    \caption{Consolidated learning curves for both cage configurations. Rotating Cage (left) and Fixed Cage (right). Each panel shows training $R^2$ and validation $R^2$ against training-set size with shared vertical axis limits, enabling direct comparison of the overfitting margin between the two cage types.}
    \label{fig:rf_learning_consolidated}
\end{figure}

The learning curves for both cage configurations (Figure~\ref{fig:rf_learning_consolidated}) exhibit a clear overfitting signature: a persistent gap of approximately 0.10 between the training $R^2$ (around 0.80) and the validation $R^2$ (around 0.70). The training score indicates how well the model has memorised the data it was shown; the validation score measures how well it generalises to held-out examples. A gap this wide means the model is capturing noise patterns in the training set that do not transfer to new flights. While the validation $R^2$ of $\sim$0.70 confirms meaningful predictive power, the size of the training--validation gap implies that performance would benefit either from acquiring additional training examples, from stronger regularisation that constrains the model's capacity to memorise, or from reducing model complexity to better match the effective information content of the 14-feature input space.

\subsubsection{Model Comparison Against Linear Baseline}

The Random Forest succeeded in recovering the impact angle from IMU data alone. But a critical sanity check remains: could a simple linear regression on the single best-correlated feature have achieved the same result? To answer this, the model was benchmarked against a Huber linear regressor operating on the single best-correlated IMU feature for each configuration.

To establish this baseline for comparison, Huber Regression was applied to the single most highly correlated IMU feature for each cage configuration (visualised in Figure~\ref{fig:rf_huber_consolidated}). Huber Regression is a robust linear technique that applies a standard squared loss function to small errors but transitions to a linear loss function for larger outliers. This makes it significantly less sensitive to the extreme, unpredictable noise spikes inherent in crash data compared to standard Ordinary Least Squares (OLS) regression. It is utilized here to determine exactly how much of the impact angle can be explained by a simple, direct linear relationship before justifying the complexity of the Random Forest.

\begin{figure}[htbp]
    \centering
    \includegraphics[width=0.85\textwidth]{plots/consolidated_top_features.png}
    \caption{Scatter plots of the top three IMU features (by Pearson~$|r|$) against impact angle, with Huber robust trendlines superimposed. These represent the best-case linear scenario: the single strongest individual sensor channel available. The wide scatter and flat slopes, particularly for the Rotating Cage, visually confirm why single-feature linear regression cannot recover the impact angle.}
    \label{fig:eda_top3_scatter}
\end{figure}

\begin{figure}[htbp]
    \centering
    \includegraphics[width=0.85\textwidth]{plots/consolidated_model_comparison.png}
    \caption{Model comparison: Random Forest vs.\ Huber linear baseline for Rotating Cage (left) and Fixed Cage (right). Bar charts show $R^2$ with $\pm1\sigma$ error bars across folds, with identical vertical axis limits for direct comparison.}
    \label{fig:rf_huber_consolidated}
\end{figure}

The comparison reveals a striking asymmetry between the two cage types. For the Fixed Cage, the single-feature Huber baseline achieves an $R^2$ of 0.348 --- weak but containing genuine correlation. The Random Forest lifts this to $R^2 = 0.731$, a gain of $\Delta R^2 = +0.384$ attributable to its ability to combine non-linear interactions across all 14 features.

For the Rotating Cage, the pattern is fundamentally different. The Huber baseline collapses to an $R^2$ of 0.004 --- functionally zero, indicating that no single linear IMU channel carries usable angle information. Yet the Random Forest recovers an $R^2$ of 0.719 from the same data. The $\Delta R^2 = +0.714$ gain is nearly double the Fixed Cage improvement, demonstrating that the rotating bearing mechanism disperses the angle signal across a broader, more distributed fingerprint that only a multi-feature non-linear model can decode. This finding is consistent with the mechanical interpretation: by decoupling the contact surface from the airframe, the rotating cage spreads collision energy across a richer set of vibrational modes, making the angle information more complex to extract but not fundamentally absent.

\subsubsection{Cross-Condition Transfer: Are the Dynamics Transferable?}

A critical practical question is whether a model trained on one cage type can perform on the other --- i.e., whether the learned vibration-to-angle mapping generalises across mechanical configurations.

\begin{figure}[htbp]
    \centering
    \includegraphics[width=\textwidth]{plots/consolidated_cross_condition_transfer.png}
    \caption{Consolidated cross-condition transfer: Fixed Cage model evaluated on Rotating Cage data (left) and Rotating Cage model evaluated on Fixed Cage data (right). Grey points show within-condition CV performance; red diamonds show the model's predictions when applied directly to opposite-condition flights. The systematic off-diagonal displacement in both directions confirms that the learned vibration-to-angle mapping does not transfer across cage mechanics.}
    \label{fig:rf_transfer_consolidated}
\end{figure}

The cross-condition transfer yields poor predictive performance in both directions, indicating that the learned models do not generalize across the two mechanical configurations. The transfer results are informative about the underlying mechanical divergence. The Fixed Cage model, when applied to Rotating Cage data, produces an $R^2$ of $-0.095$ --- worse than predicting the mean. The Rotating Cage model on Fixed Cage data fares marginally better at $R^2 = 0.257$, but still falls far below the within-condition performance of either model. A systematic bias is visible in both transfers: the predictions follow a distinct off-diagonal band, suggesting that a linear correction could partially align the two domains but that the underlying feature-to-angle relationship is fundamentally altered by the cage mechanics. The bearing mechanism does not merely add noise to an otherwise stable signal --- it changes which IMU channels carry the information and how they combine, rendering a model trained under one mechanical regime invalid in the other.

\subsubsection{Feature Importance: Which IMU Channels Carry the Signal?}

To understand which physical signatures the models rely on, the permutation-based feature importance was examined for both configurations (Figure~\ref{fig:consolidated_feature_importance}). Permutation importance measures the drop in $R^2$ when a single feature's values are randomly shuffled, breaking its association with the target while preserving its marginal distribution. A larger drop indicates a more critical feature.

\begin{figure}[htbp]
    \centering
    \includegraphics[width=\textwidth]{plots/consolidated_feature_importance.png}
    \caption{Permutation feature importance for the Random Forest models trained on Rotating Cage (left) and Fixed Cage (right). Features are ranked by the mean drop in $R^2$ when their values are permuted; error bars show $\pm 1\sigma$ across 30 permutation repetitions. The differing rank orders between cage types reflect the mechanical redistribution of collision energy by the rotating bearing.}
    \label{fig:consolidated_feature_importance}
\end{figure}

Across both configurations, vibration-derived features in the gyroscope domain dominate the importance rankings, consistent with the physical expectation that angular rate perturbations carry richer information about the impact direction than linear accelerations alone. The precise ordering differs between cage types, reflecting the mechanical redistribution of collision energy described in the Discussion chapter.

\medskip
Together, these results confirm that standard flight controller IMUs carry sufficient information for tactile angle estimation, and that the rotating cage disperses rather than destroys this signal. However, the failure of the cross-condition transfer explicitly demonstrates that this predictive mapping is highly non-linear and deeply platform-specific.
\endinput
